\section{特征函数反演与 Laplace 变换}
\label{sec:08-Numerical-Analysis/09-Complex-Analysis-and-Residues/03}

一台设备做完辨识，拿到传递函数

\[
\hat h(s)=\frac{1}{(s+1)(s^2+0.2s+25)} .
\]

要回答的问题很实际：敲它一下，十秒后还剩多少抖动？

标准答案是把 \(\hat h\) 反演回时域求 \(h(t)\)。真去做的人不多。工程上直接读极点：\(s=-1\) 对应 \(e^{-t}\)，\(s=-0.1\pm 5i\) 对应 \(e^{-0.1t}\cos(5t+\phi)\)。十秒之后，前一项衰到 \(4.5\times10^{-5}\)，后一项还有 \(0.37\)——抖动全归那对复极点，那个实极点早就不见了。谁衰得慢，看谁的实部更靠右；抖多快，看谁的虚部更大。

这个读数不是经验法则，它是围道积分的直接推论。而同一句话在概率论里叫特征函数反演，在常微分方程里叫特征根，在控制里叫稳定性判据。本节把这几件事按到同一个平面上。

\subsection{反演公式是一条积分路径}

概率论里的特征函数 \(\varphi_X(t)=\E[e^{itX}]\) 就是密度的 Fourier 变换。若 \(X\) 有密度 \(f\) 且 \(\varphi\) 可积，反演公式给出

\[
f(x)=\frac1{2\pi}\int_{-\infty}^{\infty}e^{-itx}\varphi(t)\,dt .
\]

离散情形换成 Fourier 级数的系数公式，\(p_k=\frac1{2\pi}\int_{-\pi}^{\pi}e^{-itk}\varphi(t)\,dt\)。\(\varphi\) 可积这个前提对应 \(f\) 有一定光滑性；\(\varphi\) 不可积时密度未必存在（离散分布、奇异分布都是例子），得改用 Lévy 反演定理，先算分布函数在两点之差，再取截断极限。

矩母函数 \(M_X(s)=\E[e^{sX}]\) 与特征函数是同一个变换在复平面上不同直线的取值：\(\varphi\) 取虚轴，\(M\) 取实轴。虚轴上 \(\abs{e^{itX}}=1\)，积分总收敛，所以特征函数对任何分布都存在；实轴上 \(e^{sX}\) 可能爆掉，于是 \(M\) 只在一条竖直带 \(s_-<\operatorname{Re}s<s_+\) 内有定义，带宽由分布尾部的轻重决定。\hyperref[sec:05-Probability/07-Transforms-and-Bounds/01]{``生成函数、矩母函数与特征函数''}把这两个变换并排讲过；从复平面上看，它们是一张图上的两条线，而尾部越重，收敛带越窄，最近的奇点离虚轴越近。

\subsection{看到反演积分，先检查解析性}

Cauchy 分布是个好教训。它的特征函数是 \(\varphi(t)=e^{-\abs t}\)，反演公式给出

\[
f(x)=\frac1{2\pi}\int_{-\infty}^{\infty}e^{-itx}e^{-\abs t}\,dt .
\]

看到 Fourier 积分就想封围道，这里会当场出岔子：\(e^{-\abs t}\) 在 \(t=0\) 有折角，没有跨越原点的解析延拓，任何绕过原点的围道都无从谈起。正确做法是在 \(t=0\) 处把正负半轴分开，各自是一个能收敛的初等积分：

\[
\int_{-\infty}^{\infty}e^{-itx}e^{-\abs t}\,dt
=\int_0^\infty e^{-(1+ix)t}\,dt+\int_0^\infty e^{-(1-ix)t}\,dt
=\frac{2}{1+x^2},
\]

于是 \(f(x)=1/(\pi(1+x^2))\)，Cauchy 密度。

这个结果与上一节严丝合缝。\hyperref[sec:08-Numerical-Analysis/09-Complex-Analysis-and-Residues/02]{``留数定理计算实积分''}算过 \(\int\cos(ax)/(1+x^2)\,dx=\pi e^{-\abs a}\)，那正是同一对 Fourier 关系的另一个方向；那里的绝对值来自``往哪半平面封''的二选一，这里的折角来自同一件事——\(\pm i\) 那对极点分居实轴两侧，谁被圈进去取决于指数因子往哪衰减。折角与半平面切换是同一个现象在两端的读数。

Gaussian 则相反地省事：\(\varphi(t)=e^{-t^2/2}\) 是整函数，全平面无奇点，反演积分配方即可，用不上留数。它的特征函数与密度落在同一个函数族里，这一条在很多推导里被默默用掉。

所以拿到一个反演积分，第一步不是选围道，是问被积函数能不能延拓到你想画的那条路径周围。

\subsection{最右边的奇点写着长期行为}

Laplace 变换 \(\hat f(s)=\int_0^\infty e^{-sx}f(x)\,dx\) 在 \(\operatorname{Re}s>\sigma_0\) 内收敛，\(\sigma_0\) 称为收敛横坐标，由 \(f\) 的增长速度决定。反演是 Bromwich 积分

\[
f(x)=\frac1{2\pi i}\int_{c-i\infty}^{c+i\infty}e^{sx}\hat f(s)\,ds,\qquad c>\sigma_0,
\]

沿一条竖直线积。这条线必须站在所有奇点的右边，否则积分不收敛。

把这条线往左推，就是上一节的封闭围道换了个方向：用一段左半平面的大圆弧把竖直线接成闭合回路，若弧上贡献趋于零，留数定理给出

\[
f(x)=\sum_j\operatorname{Res}\bigl(e^{sx}\hat f(s),\,s_j\bigr),
\]

求和跑遍被扫过的奇点。一阶极点 \(s_j\) 贡献 \(c_j e^{s_j x}\)，\(m\) 阶极点贡献 \(x^{m-1}e^{s_j x}\) 量级的项。有理传递函数做部分分式分解再逐项查表，做的就是这件事，只是没把留数的名字叫出来。

\begin{center}
\begin{tikzpicture}[scale=0.85,font=\small,>=stealth]
  \begin{scope}
    \draw[black!30,dashed] (0,-2.2) -- (0,2.2);
    \draw[->,black!55] (-3.7,0) -- (1.9,0) node[right,font=\footnotesize] {\(\operatorname{Re}s\)};
    \draw[->,black!55] (0,-2.4) -- (0,2.6) node[above,font=\footnotesize] {\(\operatorname{Im}s\)};
    \draw[red!60!black,thick] (1.25,-2.3) -- (1.25,2.3);
    \node[font=\footnotesize,red!55!black] at (1.25,2.72) {\(\operatorname{Re}s=c\)};
    \foreach \px/\py/\col in {-2.4/0/{blue!65!black},-0.7/1.35/{orange!80!black},-0.7/-1.35/{orange!80!black},0.5/0/{red!70!black}}{
      \draw[\col,thick] (\px-0.13,\py-0.13) -- (\px+0.13,\py+0.13);
      \draw[\col,thick] (\px-0.13,\py+0.13) -- (\px+0.13,\py-0.13);
    }
    \draw[->,black!60] (1.0,-1.95) -- (-0.6,-1.95);
    \node[below,font=\footnotesize,black!70] at (0.2,-2.0) {围道左移，扫过奇点};
  \end{scope}
  \begin{scope}[shift={(6.9,1.75)}]
    \draw[black!35] (0,0) -- (3.3,0);
    \draw[blue!65!black,thick] plot[domain=0:3.3,samples=60] (\x,{0.6*exp(-2.2*\x)});
    \node[right,font=\footnotesize,blue!55!black] at (3.45,0.1) {\(e^{-2.2x}\)：很快消失};
  \end{scope}
  \begin{scope}[shift={(6.9,0)}]
    \draw[black!35] (0,0) -- (3.3,0);
    \draw[orange!80!black,thick] plot[domain=0:3.3,samples=120] (\x,{0.6*exp(-0.7*\x)*cos(4.2*\x r)});
    \node[right,font=\footnotesize,orange!70!black] at (3.45,0.1) {\(e^{-0.7x}\cos(4.2x)\)：衰减振荡};
  \end{scope}
  \begin{scope}[shift={(6.9,-1.75)}]
    \draw[black!35] (0,0) -- (3.3,0);
    \draw[red!70!black,thick] plot[domain=0:3.3,samples=60] (\x,{0.16*exp(0.55*\x)});
    \node[right,font=\footnotesize,red!60!black] at (3.45,0.35) {\(e^{0.55x}\)：发散};
  \end{scope}
\end{tikzpicture}

\emph{每个极点在时域里换来一项 \(e^{s_jx}\)：实部定衰减率，虚部定振荡频率。竖直虚线是生死线，最右那个极点站在它哪一侧，决定整个解的命运。}
\end{center}

图里的规律可以一句话说完：奇点的实部给衰减率，虚部给振荡频率，而所有项当中衰得最慢的那一项——也就是实部最大、位置最靠右的那个奇点——统治了长期行为。开头那台设备的读数，就是这张图的一次具体填空。

这个判据在别处已经出现过好几次，只是名字不同。\hyperref[sec:15-Differential-Equations-and-Dynamical-Systems/01-ODE-Theory/02]{``线性系统与矩阵指数''}里，\(\dot y=Ay\) 的解由 \(e^{At}\) 给出，长期行为看 \(A\) 的谱横坐标（特征值实部的最大值）；\hyperref[sec:08-Numerical-Analysis/06-ODE/03]{``稳定区域与刚性方程''}里，数值格式稳定与否看 \(\lambda h\) 落在稳定区域的哪一侧，而 \(\lambda\) 正是那些特征值。三者说的是同一件事：对常系数线性方程做 Laplace 变换，微分变成乘 \(s\)，方程变成代数式，分母恰好是特征多项式，于是 \(\hat f\) 的极点就是特征根。特征根、矩阵特征值、传递函数极点——三个词，一处位置。

沿这条线还能读出刚性的来源。刚性方程的特征值实部相差几个数量级，翻译成极点，就是有的极点远在左边（对应飞快衰掉的分量），有的贴着虚轴（对应慢慢演化的分量）。显式格式的步长被最左边那个极点卡住，而人们要看的是最右边那个。刻度差得越远，格式越难做。

\subsection{生成函数的奇点决定计数序列的形状}

同一条原理在离散侧也成立。生成函数 \(G(z)=\sum_k p_kz^k\) 在 \(\abs z<R\) 内解析，\(R\) 是最近奇点到原点的距离，于是 \(p_k\) 的指数增长率立刻是 \(R^{-k}\)——这只是幂级数收敛半径的另一种说法。更细的形状由奇点的类型决定：一阶极点给出干净的 \(p_k\sim CR^{-k}\)；代数奇点 \((1-z/R)^{-\alpha}\) 在指数因子上乘一个多项式修正 \(k^{\alpha-1}\)；对数奇点带来 \(1/\log k\) 一类的慢因子。取出这些修正的标准手法是把求和圈用 Hankel 型围道贴着奇点绕过去，那是留数定理在一类退化情形下的推广。

分析组合学整门学科就架在这上面：分枝过程的灭绝概率、随机树的高度分布、组合结构的渐近计数，都先找生成函数的主奇点，再读它的类型。与 Laplace 那一侧的区别只是坐标——那边看``最右''，这边看``最近''，做的是同一个动作。

\subsection{知道位置，不等于算得出数值}

奇点分析给的是渐近形状，不是任意精度的数值。Bromwich 积分在数值上出了名地难：被积函数含 \(e^{sx}\)，沿竖直线积分时振荡剧烈，直接上求积公式会被抵消吃掉有效位——那正是\hyperref[sec:08-Numerical-Analysis/01-Floating-Point-and-Error/02]{``误差如何产生与放大''}里的老问题。可用的数值反演方法（Talbot 围道、Weeks 展开、Euler 加速）几乎都在做同一件事：把积分路径变形成一条能让 \(e^{sx}\) 衰减的曲线，让被积函数在路径上迅速趋零，从而少取几个点就够。路径变形之所以合法，仍然是 Cauchy 定理。

概率论那一侧的限度是另一种。特征函数唯一确定分布，Lévy 连续性定理还保证特征函数逐点收敛蕴含分布弱收敛——\hyperref[sec:05-Probability/08-Limit-Theorems/04]{``特征函数怎样证明中心极限定理''}整篇建立在这上面。但从 \(\varphi\) 逼近到密度逼近的这一步会掉精度：\(\varphi\) 上的小扰动可以让反演出的密度面目全非，因为反演是一次除法式的放大。可用的是分布层面的结论，不可用的是逐点密度的数值。

整个单元只用了一个动作：把积分路径在复平面上挪动。挪而不碰奇点，得到系数的指数衰减；挪到把奇点圈住，得到实积分的闭式；挪到扫过奇点，得到解的渐近展开。实轴上看到的衰减速度、振荡频率、稳定与否，都是同一批奇点在不同投影下的影子。往下的路分两条：\hyperref[sec:08-Numerical-Analysis/08-Fourier-and-Spectral-Methods/04]{``谱方法：用光滑性换高精度''}是把这套判断用在数值格式的选型上，而 Part 15 的动力系统会把``最右奇点''换成非线性场景下的不变流形与 Lyapunov 指数——那时线性叠加不再成立，位置也不再由几个孤立的点写完。
